\documentclass[11pt,a4paper]{article}
\usepackage[utf8]{inputenc}
\usepackage{amsmath}
\usepackage{graphicx}
\usepackage{booktabs}
\usepackage{hyperref}
\usepackage{listings}
\usepackage{xcolor}
\usepackage{geometry}
\usepackage{subcaption}
\usepackage{float}

\geometry{margin=1cm}

\title{\textbf{CNN Training Report} \\ 
Image Classification with Convolutional Neural Networks}
\author{Avinash C (23B1064) \and M V Pranay (23B1073)}
\date{\today}

\begin{document}

\maketitle

\section{Introduction}
This report documents the design, implementation, and evaluation of a Convolutional Neural Network (CNN) for image classification tasks. Two models were developed: Model 1 for 10-class digit recognition (Dataset 1) and Model 2 for 100-class object recognition (Dataset 2).

\section{Model Architecture Design and Rationale}

\subsection{Model 1: 10-Class Digit Recognition}

\subsubsection{Architecture}
\begin{itemize}
    \item \textbf{Input Layer:} Grayscale images of size $1 \times 32 \times 32$
    \item \textbf{Conv1:} $1 \rightarrow 8$ channels, $3 \times 3$ kernel, stride=1, padding=1
    \item \textbf{ReLU Activation}
    \item \textbf{MaxPool1:} $2 \times 2$ pool, stride=2 ($32 \times 32 \rightarrow 16 \times 16$)
    \item \textbf{Flatten:} $[B, 8, 16, 16] \rightarrow [B, 2048]$
    \item \textbf{FC1:} $2048 \rightarrow 10$ (output logits)
\end{itemize}

\subsubsection{Design Rationale}
\begin{itemize}
    \item \textbf{Single Convolutional Layer:} Sufficient for simple digit recognition task
    \item \textbf{8 Feature Maps:} Adequate capacity to capture basic edges and patterns in digits
    \item \textbf{3×3 Kernel:} Standard kernel size balancing receptive field and computational cost
    \item \textbf{Padding=1:} Maintains spatial dimensions after convolution
    \item \textbf{MaxPooling:} Provides spatial invariance and reduces computation
\end{itemize}

\subsection{Model 2: 100-Class Object Recognition}
\subsubsection{Architecture}
\begin{itemize}
    \item \textbf{Input Layer:} Grayscale images of size $1 \times 32 \times 32$
    \item \textbf{Block 1:}
        \begin{itemize}
            \item \textbf{Conv1:} $1 \rightarrow 8$ channels, $3 \times 3$ kernel, stride=1, padding=1
            \item \textbf{ReLU and MaxPool:} $2 \times 2$ pool ($32 \times 32 \rightarrow 16 \times 16$)
        \end{itemize}
    \item \textbf{Block 2:}
        \begin{itemize}
            \item \textbf{Conv2:} $8 \rightarrow 16$ channels, $3 \times 3$ kernel, stride=1, padding=1
            \item \textbf{ReLU and MaxPool:} $2 \times 2$ pool ($16 \times 16 \rightarrow 8 \times 8$)
        \end{itemize}
    \item \textbf{Block 3:}
        \begin{itemize}
            \item \textbf{Conv3:} $16 \rightarrow 32$ channels, $3 \times 3$ kernel, stride=1, padding=1
            \item \textbf{ReLU and MaxPool:} $2 \times 2$ pool ($8 \times 8 \rightarrow 4 \times 4$)
        \end{itemize}
    \item \textbf{Block 4:}
        \begin{itemize}
            \item \textbf{Conv4:} $32 \rightarrow 64$ channels, $3 \times 3$ kernel, stride=1, padding=1
            \item \textbf{ReLU and MaxPool:} $2 \times 2$ pool ($4 \times 4 \rightarrow 2 \times 2$)
        \end{itemize}
    \item \textbf{Flatten:} $[B, 64, 2, 2] \rightarrow [B, 256]$
    \item \textbf{FC1:} $256 \rightarrow 100$ (output logits)
\end{itemize}

\subsubsection{Design Rationale}
    \begin{itemize}
        \item \textbf{Four-Block Feature Funnel:} Uses a deeper architecture to extract highly abstract features necessary 
        for 100-class discrimination.
        \item \textbf{Tapered Channel Expansion:} Starts with a low channel count (8) and doubles it at every spatial 
        reduction ($32 \to 16 \to 8 \to 4 \to 2$) to maintain a balanced computational budget while increasing representational depth.
        \item \textbf{Aggressive Downsampling:} Four pooling layers reduce the final feature map to $2 \times 2$, effectively 
        aggregating global spatial information before the fully connected layer.
    \end{itemize}

\section{Trainable Parameters}

\subsection{Model 1: Parameter Count}

\textbf{Conv1 Parameters:}
\begin{align*}
\text{Weights} &= C_{out} \times C_{in} \times K_h \times K_w = 8 \times 1 \times 3 \times 3 = 72 \\
\text{Biases} &= C_{out} = 8 \\
\text{Total Conv1} &= 72 + 8 = 80
\end{align*}

\textbf{FC1 Parameters:}
\begin{align*}
\text{Weights} &= N_{in} \times N_{out} = 2048 \times 10 = 20,480 \\
\text{Biases} &= N_{out} = 10 \\
\text{Total FC1} &= 20,480 + 10 = 20,490
\end{align*}

\textbf{Total Parameters (Model 1):} $80 + 20,490 = \boxed{20,570}$

\section{Model 2: Detailed Analysis}

\subsection{Parameter Count}

\textbf{Conv1 Parameters ($1 \to 8$):}
\begin{align*}
\text{Weights} &= 8 \times 1 \times 3 \times 3 = 72 \\
\text{Biases} &= 8 \\
\text{Total Conv1} &= 80
\end{align*}

\textbf{Conv2 Parameters ($8 \to 16$):}
\begin{align*}
\text{Weights} &= 16 \times 8 \times 3 \times 3 = 1,152 \\
\text{Biases} &= 16 \\
\text{Total Conv2} &= 1,168
\end{align*}

\textbf{Conv3 Parameters ($16 \to 32$):}
\begin{align*}
\text{Weights} &= 32 \times 16 \times 3 \times 3 = 4,608 \\
\text{Biases} &= 32 \\
\text{Total Conv3} &= 4,640
\end{align*}

\textbf{Conv4 Parameters ($32 \to 64$):}
\begin{align*}
\text{Weights} &= 64 \times 32 \times 3 \times 3 = 18,432 \\
\text{Biases} &= 64 \\
\text{Total Conv4} &= 18,496
\end{align*}

\textbf{FC1 Parameters ($256 \to 100$):}
\begin{align*}
\text{Weights} &= 256 \times 100 = 25,600 \\
\text{Biases} &= 100 \\
\text{Total FC1} &= 25,700
\end{align*}

\textbf{Total Parameters (Model 2):} $80 + 1,168 + 4,640 + 18,496 + 25,700 = \boxed{50,084}$

\section{MACs and FLOPs Analysis}

\subsection{Computational Complexity Definitions}
\begin{itemize}
    \item \textbf{MAC (Multiply-Accumulate):} One multiplication followed by one addition
    \item \textbf{FLOP (Floating Point Operation):} Each multiply and add counted separately
    \item \textbf{Relationship:} $\text{FLOPs} = 2 \times \text{MACs}$ (assuming each MAC = 2 FLOPs)
\end{itemize}

\subsection{Conv2D MACs Formula}
For a convolution layer:
\begin{equation}
\text{MACs} = H_{out} \times W_{out} \times C_{out} \times (K_h \times K_w \times C_{in})
\end{equation}

\subsection{Linear Layer MACs Formula}
For a fully connected layer:
\begin{equation}
\text{MACs} = N_{in} \times N_{out}
\end{equation}

\subsection{Model 1: MACs/FLOPs per Forward Pass}

\textbf{Conv1:}
\begin{align*}
\text{MACs} &= 32 \times 32 \times 8 \times (3 \times 3 \times 1) \\
&= 1024 \times 8 \times 9 = 73,728
\end{align*}

\textbf{FC1:}
\begin{align*}
\text{MACs} &= 2048 \times 10 = 20,480
\end{align*}

\textbf{Total Model 1:}
\begin{align*}
\text{Total MACs} &= 73,728 + 20,480 = \boxed{94,208} \\
\text{Total FLOPs} &= 2 \times 94,208 = \boxed{188,416}
\end{align*}

\subsection{Model 2: MACs/FLOPs per Forward Pass}

\textbf{Conv1 ($32 \times 32 \times 8$):}
\begin{align*}
\text{MACs} &= 32 \times 32 \times 8 \times (3 \times 3 \times 1) \\
&= 1024 \times 8 \times 9 = 73,728
\end{align*}

\textbf{Conv2 ($16 \times 16 \times 16$):}
\begin{align*}
\text{MACs} &= 16 \times 16 \times 16 \times (3 \times 3 \times 8) \\
&= 256 \times 16 \times 72 = 294,912
\end{align*}

\textbf{Conv3 ($8 \times 8 \times 32$):}
\begin{align*}
\text{MACs} &= 8 \times 8 \times 32 \times (3 \times 3 \times 16) \\
&= 64 \times 32 \times 144 = 294,912
\end{align*}

\textbf{Conv4 ($4 \times 4 \times 64$):}
\begin{align*}
\text{MACs} &= 4 \times 4 \times 64 \times (3 \times 3 \times 32) \\
&= 16 \times 64 \times 288 = 294,912
\end{align*}

\textbf{FC1 ($256 \to 100$):}
\begin{align*}
\text{MACs} &= 256 \times 100 = 25,600
\end{align*}

\textbf{Total Model 2 Complexity:}
\begin{align*}
\text{Total MACs} &= 73,728 + 294,912 + 294,912 + 294,912 + 25,600 = \boxed{984,064} \\
\text{Total FLOPs} &= 2 \times 984,064 = \boxed{1,968,128}
\end{align*}

\section{Dataset Loading Time}

\subsection{Dataset Statistics}
\begin{itemize}
    \item \textbf{Dataset 1:} 10 classes (digits 0-9), grayscale images resized to $32 \times 32$
    \item \textbf{Dataset 2:} 100 classes (CIFAR-100 categories), grayscale images resized to $32 \times 32$
    \item \textbf{Preprocessing:} Normalization to [0, 1], data augmentation (random flips, rotation)
\end{itemize}

\subsection{Loading Time Measurements}
\begin{table}[H]
\centering
\begin{tabular}{lrr}
\toprule
\textbf{Metric} & \textbf{Dataset 1} & \textbf{Dataset 2} \\
\midrule
Training samples & 3.7358 & 4.5397s \\
Test samples & 3.2589 & 3.7273 \\
Batch size & 32 & 32 \\
\bottomrule
\end{tabular}
\caption{Dataset loading time measurements}
\end{table}

\section{Training and Validation Performance}

\subsection{Training Configuration}
\begin{itemize}
    \item \textbf{Optimizer:} Stochastic Gradient Descent (SGD)
    \item \textbf{Learning Rate:} 0.01
    \item \textbf{Batch Size:} 32 (training), 64 (evaluation)
    \item \textbf{Number of Epochs:} 10
    \item \textbf{Loss Function:} Cross-Entropy Loss
    \item \textbf{Data Augmentation:} Random horizontal flip, random rotation ($\pm 10^\circ$)
\end{itemize}

\subsection{Model 1: Training Performance}
\begin{table}[H]
\centering
\begin{tabular}{crr}
\toprule
\textbf{Epoch} & \textbf{Train Loss} & \textbf{Time (s)} \\
\midrule
1 & 1.0697 & 20.38 \\
2 & 0.5873 & 20.12 \\
3 & 0.5271 & 21.15 \\
4 & 0.4729 & 20.92 \\
5 & 0.4218 & 22.04 \\
6 & 0.3750 & 21.46 \\
7 & 0.3359 & 38.95 \\
8 & 0.3033 & 25.66 \\
9 & 0.2765 & 21.41 \\
10 & 0.2547 & 21.16 \\
\bottomrule
\end{tabular}
\caption{Model 1 training metrics across epochs}
\end{table}

\textbf{Final Test Accuracy:} 92.97\% (55785/60000)

\subsection{Model 2: Training Performance}
% Replace with your actual results
\begin{table}[H]
\centering
\begin{tabular}{crr}
\toprule
\textbf{Epoch} & \textbf{Train Loss} & \textbf{Time (s)} \\
\midrule
1 & 4.5374 & 129.08 \\
2 & 4.1907 & 120.63 \\
3 & 3.9313 & 130.99 \\
4 & 3.7337 & 151.13 \\
5 & 3.5913 & 138.09 \\
6 & 3.4682 & 115.92 \\
7 & 3.3625 & 118.03 \\
8 & 3.2729 & 118.07 \\
9 & 3.1994 & 119.53 \\
10 & 3.1331 & 120.73 \\
\bottomrule
\end{tabular}
\caption{Model 2 training metrics across epochs}
\end{table}

\textbf{Final Test Accuracy:} 27.53\% (13765/50000)

\section{Performance Visualization}


\begin{figure}[H]
    \centering
    \includegraphics[width=0.7\textwidth]{loss_plot_model1.png}
    \caption{Model 1: Training loss over epochs}
    \label{fig:loss1}
\end{figure}

\begin{figure}[H]
    \centering
    \includegraphics[width=0.7\textwidth]{loss_plot_model2.png}
    \caption{Model 2: Training loss over epochs}
    \label{fig:loss2}
\end{figure}

\section{Failed Design Decision}

\subsection{Attempted Design: Three Convolutional Layers}

\subsubsection{Initial Architecture}
I initially attempted a deeper architecture for Model 2 with three convolutional layers:
\begin{itemize}
    \item Conv1: $1 \rightarrow 16$ channels, MaxPool ($32 \rightarrow 16$)
    \item Conv2: $16 \rightarrow 32$ channels, MaxPool ($16 \rightarrow 8$)
    \item Conv3: $32 \rightarrow 64$ channels, MaxPool ($8 \rightarrow 4$)
    \item FC: $64 \times 4 \times 4 = 1024 \rightarrow 100$
\end{itemize}

\subsubsection{Why It Failed}

\textbf{1. Vanishing Spatial Information}
\begin{itemize}
    \item After three maxpool operations, spatial dimensions reduced to $4 \times 4$
    \item Original $32 \times 32$ images already small; excessive downsampling lost critical details
    \item Fine-grained features needed for discriminating 100 classes were destroyed
\end{itemize}

\textbf{2. Overfitting}
\begin{itemize}
    \item Training loss decreased rapidly but test accuracy remained poor 
    \item High capacity but insufficient training data
    \item Model memorized training samples rather than learning generalizable features
\end{itemize}

\end{document}
